\clearpage\thispagestyle{empty}
\pdfbookmark[0]{Quellen zum Bericht über Vesta vom September 1807}{vesta1807september-sources}
\begin{center}\textbf{Quellen zum Bericht über Vesta vom September 1807}\end{center}
\noindent Diese KI-Übertragung in heutiges Deutsch folgt Franz Xaver von Zachs Bericht in \href{https://gdz.sub.uni-goettingen.de/id/PPN236007467}{Gauß’ \textit{Werken}, Band~VI} (Göttingen, 1874), S.~290--292. Er enthält Äußerungen von Heinrich Wilhelm Matthias Olbers und Carl Friedrich Gauß sowie die Pariser Beobachtungen, die Bouvard Gauß mitteilte. Der Erstdruck erschien als Beitrag XXII, \textit{Fortgesetzte Nachrichten über den neuen von Dr. Olbers in Bremen entdeckten Planeten Vesta}, in der \textit{Monatlichen Correspondenz zur Beförderung der Erd- und Himmelskunde}, Band~XVI, September 1807, \href{https://www.digitale-sammlungen.de/en/view/bsb10538607?page=299}{S.~285}--\href{https://www.digitale-sammlungen.de/en/view/bsb10538607?page=305}{291}. Olbers’ Brief beginnt auf \href{https://www.digitale-sammlungen.de/en/view/bsb10538607?page=300}{S.~286}, die dritten Elemente stehen auf \href{https://www.digitale-sammlungen.de/en/view/bsb10538607?page=301}{S.~287} und die Vergleichstafel und Gauß’ Bemerkungen zu Hardings Atlas auf \href{https://www.digitale-sammlungen.de/en/view/bsb10538607?page=302}{S.~288}.
\par\medskip
Die Überschrift der Vergleichstafel nennt in der Werkausgabe die mittlere Zeit in Paris, im Erstdruck die mittlere Zeit in Bremen. Die beiden editierbaren Vesta-Zeichen folgen den Konturen der Werkausgabe. In der Zeitangabe für den 13. April zeigt das Göttinger Digitalisat einen hochgestellten schrägen Strich zwischen 6 und 2. Das \href{https://archive.org/details/werkecarlf06gausrich}{Exemplar der University of California Libraries} zeigt eine weitere, versetzte Spur. Die Lesung des Zeichens bleibt unsicher. Der \href{https://www.digitale-sammlungen.de/en/view/bsb10538607?page=302}{Erstdruck} gibt den Wert mit 6,2 Sekunden an. Die Werkausgabe beschreibt die Ephemeride von \href{https://www.digitale-sammlungen.de/en/view/bsb10538607?page=303}{S.~289}, \href{https://www.digitale-sammlungen.de/en/view/bsb10538607?page=304}{S.~290} und \href{https://www.digitale-sammlungen.de/en/view/bsb10538607?page=305}{S.~291} des Erstdrucks in eckigen Klammern. Die Zeitschriftendigitalisate stammen von der Bayerischen Staatsbibliothek. Übertragung und Satz wurden mit KI erstellt.
\clearpage
